% T30 source: Примеры базовой модели.tex lines 489--627
\detailedexampleheading{Variant 2. Interactive and Delegated Tasks}

In this variant, two agent streams are available ($P = 2$). Tasks that require
architectural decisions are performed interactively, whereas more formalized
tasks are delegated to an agent and accepted by the developer after an
autonomous phase. The operating mode applies to a task rather than being fixed
to a particular agent: any available agent may receive the next ready task in
either of the prescribed modes.

Tasks are assigned to operating modes as follows:
\begin{itemize}
  \item \textbf{Interactive mode}: API contract, business logic, and background
  worker---tasks that require context and architectural decisions.
  \item \textbf{Delegated mode}: persistence layer, observability, tests, and
  packaging---tasks with more formalized requirements.
\end{itemize}

The configuration is $\sigma_2 = (P{=}2,\ \rho)$ with the mixed mode assignment
$\rho$: tasks in Stream 1 are performed in interactive mode, and tasks in
Stream 2 are performed in delegated mode.

\paragraph{Parameters.}
\begin{table}[ht!]
\centering
\caption{Parameters of Variant 2}
\label{tab:practical-v2-params}
\begin{tabular}{clcccc}
\toprule
$i$ & Task & $Z_i$ & $k_i$ & $h_i$ & Mode \\
\midrule
1 & API contract and request validation & 2 & 0.75 & 0.75 & interact. \\
2 & Persistence layer & 3 & 0.75 & 0.20 & deleg. \\
3 & Business logic & 4 & 0.85 & 0.85 & interact. \\
4 & Background worker & 3 & 0.90 & 0.90 & interact. \\
5 & Observability and configuration & 2 & 0.70 & 0.20 & deleg. \\
6 & Tests & 3 & 0.70 & 0.20 & deleg. \\
7 & Packaging and deployment & 2 & 0.80 & 0.20 & deleg. \\
\bottomrule
\end{tabular}
\end{table}

In delegated mode, the coefficients $k_i$ are slightly lower than in
Variant~1. This difference results from the change in operating mode: the
agent works independently and does not wait for intermediate responses from
the developer. This effect arises even in a single stream and is consistent
with the definition of $k_i^{(r)}$ as a characteristic of the ``task
$\times$ tool $\times$ operating mode'' triple; the tool is held fixed in the
scenario. For interactive tasks, the coefficients are close to those in
Variant~1, and $h_i = k_i$ because the developer is occupied throughout their
execution. For delegated tasks, $h_i = 0.20$ is assumed; this value accounts
for task specification and review of the result. In the schedule, the human
time is divided equally between initial task specification and final review,
and the agent slot is retained over the entire interval from the beginning of
specification through completion of review.

\paragraph{Calculations.}
\begin{align*}
  \sum_{i=1}^{7} Z_i k_i &= 2{\cdot}0.75 + 3{\cdot}0.75 + 4{\cdot}0.85 + 3{\cdot}0.90 + 2{\cdot}0.70 + 3{\cdot}0.70 + 2{\cdot}0.80 \\
  &= 1.50 + 2.25 + 3.40 + 2.70 + 1.40 + 2.10 + 1.60 = 14.95.
\end{align*}

The weighted-average normalized duration is
\[
  \bar{k}_w = \frac{14.95}{19} \approx 0.787.
\]

The task weights in the graph are $w_i = Z_i k_i X$:
\[
  \{w_i\}_{i=1}^{7} = \{6.0,\; 9.0,\; 13.6,\; 10.8,\; 5.6,\; 8.4,\; 6.4\}~\text{h}.
\]

The total work is
\[
  W = \sum_{i=1}^{7} w_i = 14.95 \cdot 4 = 59.8~\text{h}.
\]

The idealized completion time for independent tasks, that is, the lower bound
based on total work when $P = 2$, is
\[
  T_{ai}^{(0)} = \frac{W}{P} = \frac{59.8}{2} = 29.9~\text{h}.
\]

The idealized speedup without accounting for dependencies is
\[
  S_E = \frac{76}{29.9} \approx 2.54.
\]

The critical path through the graph comprises Tasks
$1 \to 2 \to 3 \to 4 \to 6$:
\[
  L = 6.0 + 9.0 + 13.6 + 10.8 + 8.4 = 47.8~\text{h}.
\]

The human-attention time at level M4 is
\[
  \sum_{i=1}^{7} Z_i h_i = 2{\cdot}0.75 + 3{\cdot}0.20 + 4{\cdot}0.85 + 3{\cdot}0.90 + 2{\cdot}0.20 + 3{\cdot}0.20 + 2{\cdot}0.20 = 9.6,
\]
\[
  H = 9.6 \cdot 4 = 38.4~\text{h},
  \qquad
  \bar{h}_w = \frac{9.6}{19} \approx 0.505,
  \qquad
  \frac{1}{\bar{h}_w} \approx 1.98.
\]

The lower bound on the makespan accounting for the dependency graph and the
human resource is
\[
  B_4 = \max\left\{\frac{W}{P},\; L,\; H\right\} = \max\{29.9,\; 47.8,\; 38.4\} = 47.8~\text{h}.
\]

The assignment to streams from Table~\ref{tab:practical-v2-params} that
respects the dependencies is
\begin{center}
\small
\begin{tabularx}{\textwidth}{cXX}
\toprule
Stream & Task intervals, h & Idle time, h \\
\midrule
1 & 1: $[0,6.0]$; 3: $[15.0,28.6]$; 4: $[28.6,39.4]$ & $[6.0,15.0]$, waiting for the persistence layer \\
2 & 2: $[6.0,15.0]$; 5: $[15.0,20.6]$; 7: $[20.6,27.0]$; 6: $[39.4,47.8]$ & $[27.0,39.4]$, waiting for dependencies \\
\bottomrule
\end{tabularx}
\end{center}

If only the two agent streams and the precedence relations are considered,
this assignment has a makespan of $47.8$ hours, numerically equal to the lower
bound. However, as the following check shows, it is not a feasible level-M4
schedule.
\[
  T_{\mathrm{agents}}=B_4=47.8~\text{h}.
\]

\paragraph{Human-Resource Check (Level M4).}
The schedule above is feasible with respect to the streams and dependencies
but \emph{infeasible with respect to the human resource}. The developer is
continuously occupied by interactive tasks over the interval
$[15.0,\,39.4]$: first by Task 3 and then by Task 4. At the same time, delegated
Tasks 5 $[15.0,\,20.6]$ and 7 $[20.6,\,27.0]$ each require $1.6$ hours of human
attention for task specification and review. These intervals overlap the
interactive work, which is impossible with a single developer. Consequently,
the schedule that is formally feasible on the agent streams does not provide
the resource capacity needed for the timely acceptance of the delegated
results.

The conflict can be eliminated by restructuring the schedule: delegated tasks
and their human-attention phases are placed in intervals when the developer is
available. Once started, a task continuously retains its assigned agent slot,
including while it waits for final review. The corrected schedule is as
follows:
\begin{itemize}
  \item Agent 1: Task 1 $[0,\,6.0]$; Task 2 $[6.0,\,15.0]$; Task 3
  $[15.0,\,28.6]$; Task 4 $[28.6,\,39.4]$; Task 6 $[39.4,\,47.8]$.
  \item Agent 2: Task 5 $[7.2,\,12.8]$; Task 7 $[12.8,\,41.4]$.
  For Task 7, specification occupies $[12.8,\,13.6]$, autonomous work occupies
  $[13.6,\,18.4]$, waiting for review occupies $[18.4,\,40.6]$, and review
  occupies $[40.6,\,41.4]$.
  \item Developer: Task 1 $[0,\,6.0]$; specification of Task 2
  $[6.0,\,7.2]$; specification of Task 5 $[7.2,\,8.0]$; review of Task 5
  $[12.0,\,12.8]$; specification of Task 7 $[12.8,\,13.6]$; review of Task 2
  $[13.8,\,15.0]$; Task 3 $[15.0,\,28.6]$; Task 4 $[28.6,\,39.4]$;
  specification of Task 6 $[39.4,\,40.6]$; review of Task 7
  $[40.6,\,41.4]$; review of Task 6 $[46.6,\,47.8]$. The developer intervals
  are pairwise nonoverlapping, and their total duration is $38.4$ hours $= H$.
\end{itemize}
All dependencies are satisfied. Task 3 starts at time $15.0$, after Task 2 has
been fully accepted; Task 7 starts after Task 5 has been accepted; and Task 6
starts after Task 4. The waiting interval of Task 7 is not used for other work
on Agent 2. Despite this local blocking, the makespan remains $47.8$ hours
because Task 7 is not on the critical path. Thus, the human resource does not
increase the makespan only when the phases and the queue are planned
explicitly; the original schedule was infeasible.

The corrected schedule is feasible and attains the lower bound; it is
therefore optimal:
\[
  T(\pi_2)=T^{*}=B_4=47.8~\text{h}.
\]

The attainable speedup accounting for the graph and the human resource is
\[
  S_E^{\text{ach}} = \frac{76}{47.8} \approx 1.59.
\]

\paragraph{Interpretation.}
If the tasks were treated as independent, the lower bound $W/P=29.9$ hours
would correspond to a speedup of approximately $2.54\times$. However, the
dependency graph substantially reduces the exploitable parallelism: Tasks
$1 \to 2 \to 3 \to 4 \to 6$ form a critical path of length $47.8$ hours.
After the graph is taken into account, the attainable speedup decreases to
$1.59\times$. Observability and deployment do proceed in parallel, but they
are outside the critical path and have almost no effect on the overall
makespan.

The key observation is that parallelism $P = 2$ helps only where the graph
permits independent work. In this variant, most of the effort lies on a
sequential chain, so balancing the streams without accounting for dependencies
overstates the effect of the configuration in the illustrative scenario.

The human resource is not yet the active constraint: $H = 38.4$ hours
$< L = 47.8$ hours, and the attained speedup remains below the ceiling
$1/\bar{h}_w \approx 1.98$ ($1.59 < 1.98$). This residual capacity can be used,
however, only when review intervals are explicitly scheduled. The three
interactive tasks with $h = k$ occupy the developer for a total of $30.4$
hours, so reviews of delegated tasks must be placed in the remaining
intervals.
